
\label{fasttransportsection} This simulation is here presented to illustrate
code capabilities for a realistic magnetic configuration and a hot deuterium
plasma. It corresponds to a typical JET Lower Hybrid current drive
discharge, where both bounce-averaging and fast electron particle transport
are considered. Global parameters of the discharge are gathered in section
``JET'' of the tokamak parameter M-file ``ptok\_dke\_1yp.m''.

As shown in Fig.\ref{radialgridJET} , the spatial $\psi $ half-grid, on
which the electron distribution function is calculated, is non-uniform,
though grid steps $\Delta \rho $ for the normalized radius defined in Sec. 
\ref{coordinatesystemsection} are constant for this example. The spatial
mesh size for $f_{0}^{\left( 0\right) }$ has $n_{\psi }-1=14$ points, which
is the typical number of spatial grid points considered in most of the
current drive simulations. The corresponding pitch-angle grid is strongly
non uniform, as indicated in Fig.\ref{pitchanglegridJET}, since
trapped/passing boundaries corresponding to all radial positions are exactly
placed on the flux grid in momentum space. The total number of pitch-angle
points for $f_{0}^{\left( 0\right) }$ is $n_{\xi _{0}}-1=206$. Finally the
momentum grid is also taken no uniform with $n_{p}-1=210$ points. As shown
in Fig.\ref{momentumgridJET} the momentum grid is first nearly uniform in
the vicinity of $p=0$ for the first twenty points, with a small step $\Delta %
p\simeq 0.07$ , as indicated in Fig. \ref{momentumgridstepJET}. For $%
p\gtrsim 1.5,$ the grid is also uniform, but with larger momentum step $%
\Delta p\simeq 0.15$. The smooth transition between the two grids involves $7
$ grid points.

%TCIMACRO{
%\TeXButton{Fig_radialgridJET}{\fig{figures/radialgridJET.eps}{Radial grid for 3-D JET current drive simulation. Circles correspond to the normalized poloidal flux coordinate $\psi $,
%while crosses correspond to normalized radius $\rho $}{radialgridJET}{7cm} }}
%BeginExpansion
\fig{figures/radialgridJET.eps}{Radial grid for 3-D JET current drive simulation. Circles correspond to the normalized poloidal flux coordinate $\psi $,
while crosses correspond to normalized radius $\rho $}{radialgridJET}{7cm} %
%EndExpansion

%TCIMACRO{
%\TeXButton{Fig_pitchanglegridJET}{\fig{figures/pitchanglegridJET.eps}{Pitch-angle grid for 3-D JET current drive simulation. }{pitchanglegridJET}{7cm} }}
%BeginExpansion
\fig{figures/pitchanglegridJET.eps}{Pitch-angle grid for 3-D JET current drive simulation. }{pitchanglegridJET}{7cm} %
%EndExpansion

%TCIMACRO{
%\TeXButton{Fig_momentumgridJET}{\fig{figures/momentumgridJET.eps}{Momentum grid for 3-D JET current drive simulation. }{momentumgridJET}{7cm} }}
%BeginExpansion
\fig{figures/momentumgridJET.eps}{Momentum grid for 3-D JET current drive simulation. }{momentumgridJET}{7cm} %
%EndExpansion

%TCIMACRO{
%\TeXButton{Fig_momentumgridstepJET}{\fig{figures/momentumgridstepJET.eps}{Momentum grid step for 3-D JET current drive simulation.  Circles correspond to the flux grid,
%while stars to the distribution function half-grid}{momentumgridstepJET}{7cm} }}
%BeginExpansion
\fig{figures/momentumgridstepJET.eps}{Momentum grid step for 3-D JET current drive simulation.  Circles correspond to the flux grid,
while stars to the distribution function half-grid}{momentumgridstepJET}{7cm} %
%EndExpansion

Temperature and density profiles used as input for the magnetic equilibrium
calculations done by HELENA code are presented in Fig. \ref{profJET}. As
usual in current drive regime at low density, the electron temperature
profile is peaked, and its value is much larger than ion temperature. The
dominant ion (deuterium) profile is determined selfconsistently from the
electron density and effective charge profiles, the latter being taken flat
at a level of $Z_{eff}=1.5$. In the calculations which are based on
electroneutrality, only one fully stripped impurity is considered, i.e.
carbon $Z_{s}=6$, in order to avoid the use of an impurity transport code.
The resulting $2-D$ contour plot of the magnetic poloidal flux surfaces is
given in Fig. \ref{equilJET}, for a configuration with a bottom X-point
close to the divertor. Here, calculations are performed for a monotonic
safety factor profile.

%TCIMACRO{
%\TeXButton{Fig_profJET}{\fig{figures/profJET.eps}{Ion and electron temperature and density profiles, and effective charge profile used for calculating the JET magnetic equilibrium with HELENA. 
%Here hydrogen and tritium densities are zero (pure deuterium plasma) . The poloidal flux coordinate $\psi $ as function of the normalized radius $\rho $ 
%in the equatorial mid-plane corresponds to the magnetic equilibrium code output}{profJET}{7cm} }}
%BeginExpansion
\fig{figures/profJET.eps}{Ion and electron temperature and density profiles, and effective charge profile used for calculating the JET magnetic equilibrium with HELENA. 
Here hydrogen and tritium densities are zero (pure deuterium plasma) . The poloidal flux coordinate $\psi $ as function of the normalized radius $\rho $ 
in the equatorial mid-plane corresponds to the magnetic equilibrium code output}{profJET}{7cm} %
%EndExpansion

%TCIMACRO{
%\TeXButton{Fig_equilJET}{\fig{figures/equilJET.eps}{$2-D$ contour plot of the poloidal magnetic flux surfaces as calculated for JET tokamak by the code HELENA}{equilJET}{7cm} }}
%BeginExpansion
\fig{figures/equilJET.eps}{$2-D$ contour plot of the poloidal magnetic flux surfaces as calculated for JET tokamak by the code HELENA}{equilJET}{7cm} %
%EndExpansion

In the calculations, fully relativistic corrections are considered. In Fig.%
\ref{vvspJET}, the deviation from the relation $v=p$ clearly indicates that
relativistic corrections come into play when $p\gtrsim 4\footnote{%
Here $v$ and $p$ are normalized respectively to $p_{th}/m_{e}$ and $p_{th}$,
where $m_{e}$ is the electron rest mass and $p_{th}$ the thermal
momentum.All reference values in the calculations are taken at the plasma
center, including  $p_{th}$ as discussed in Sec.\ref{normalization}.}.$
Consequently, since a simplified expression for the Lower Hybrid quasilinear
diffusion operator is considered, as described in Appendix \ref
{simplifiedlhcd} and used in Sec. \ref{lhcdsection}, boundary values of the
resonance domain which correspond to $\left( v_{1}=3.5,v_{2}=6\right) $
exhibit strong curvature as a function of $p$, as shown in Fig. \ref
{DlhpeakJET}. In order to simulate a spatialy localized off-axis absorption
of the RF power, the quasilinear diffusion coefficient $\overline{D}%
_{0,new}^{LH}\left( \psi ,p\right) $is defined as 
\begin{equation}
\overline{D}_{0,new}^{LH}\left( \psi ,p\right) =\overline{D}%
_{0,new}^{LH}\left( 0,p\right) \frac{2\sqrt{\ln 2}}{\sqrt{\pi }}\exp \left[ -%
\frac{\left( \rho -\rho _{LH}^{peak}\right) ^{2}}{\left( \frac{\Delta \rho
_{LH}}{2\sqrt{\ln 2}}\right) ^{2}}\right] 
\end{equation}
where $\overline{D}_{0,new}^{LH}\left( 0,p\right) =1$ in the interval $v_{1}%
\leq v_{\Vert }\leq v_{2},$ and $\rho _{LH}^{peak}=0.5$ is the radial
position at which absorption is maximum, and $\Delta \rho _{LH}=0.2$ the
half-width. For $\rho \leq 0.2$ and $\rho \geq 0.65,$ power absorption is
null. 

%TCIMACRO{
%\TeXButton{Fig_vvspJET}{\fig{figures/vvspJET.eps}{Momentum dependence of the relativistic Maxwellian distribution function at $\rho\simeq 0.36$, 
%and relation between velocity $v$ and momentum $p$. The deviation from the main diagonal indicates that above $p=4$,
%relativistic effects become important}{vvspJET}{7cm} }}
%BeginExpansion
\fig{figures/vvspJET.eps}{Momentum dependence of the relativistic Maxwellian distribution function at $\rho\simeq 0.36$, 
and relation between velocity $v$ and momentum $p$. The deviation from the main diagonal indicates that above $p=4$,
relativistic effects become important}{vvspJET}{7cm} %
%EndExpansion

%TCIMACRO{
%\TeXButton{Fig_DlhpeakJET}{\fig{figures/DlhpeakJET.eps}{$2-D$ contour plot in momentum space of the Lower Hybrid quasilinear diffusion cofficient at $\rho\simeq 0.36$.
%The relativistic curvature of the lower bound of the resonance domain avoid intersection with the region of trapped electrons. 
%The two full straight lines correspond to trapped/passing boundaries at that radial position}{DlhpeakJET}{7cm} }}
%BeginExpansion
\fig{figures/DlhpeakJET.eps}{$2-D$ contour plot in momentum space of the Lower Hybrid quasilinear diffusion cofficient at $\rho\simeq 0.36$.
The relativistic curvature of the lower bound of the resonance domain avoid intersection with the region of trapped electrons. 
The two full straight lines correspond to trapped/passing boundaries at that radial position}{DlhpeakJET}{7cm} %
%EndExpansion

Fast electron radial transport that could result from magnetic turbulence is
expressed in the usual form
\begin{equation}
D_{\psi }^{\left( 0\right) }\left( \psi ,p\right) =D_{\psi }^{\left(
0\right) }\left( 0,p\right) H\left( \left| v_{\Vert }\right| -v_{\Vert
c}\right) \left( \left| v_{\Vert }\right| -v_{\Vert c}\right) /v_{\Vert c}
\end{equation}
where $D_{\psi }^{\left( 0\right) }$ is the radial diffusion coefficient as
defined in Sec.\ref{normalization}. Here, $D_{\psi }^{\left( 0\right) }$
scales like $\left| v_{\Vert }\right| $ as expected from theory \cite{pey93}%
, with a velocity threshold which is set to $3.5$. The simulation is
performed with $D_{\psi }^{\left( 0\right) }=5m^{2}/s$, its value and the
threshold level being uniform throughout the plasma, while any convection is
neglected. The $2-D$ contour plot of the $D_{\psi }^{\left( 0\right) }$ is
shown in Fig. \ref{DradialJET}

%TCIMACRO{
%\TeXButton{Fig_DradialJET}{\fig{figures/DradialJET.eps}{On the left side, $2-D$ contour plot of the radial diffusion rate at $\rho\simeq 0.36$. The velocity threshold corresponds to a kinetic energy
%of 35 keV approximately in the $MKSA$ units. On the right side, the velocity dependence of $D_{\psi }^{\left( 0\right) }$ at
%$\xi_{0}=1$}{DradialJET}{7cm} }}
%BeginExpansion
\fig{figures/DradialJET.eps}{On the left side, $2-D$ contour plot of the radial diffusion rate at $\rho\simeq 0.36$. The velocity threshold corresponds to a kinetic energy
of 35 keV approximately in the $MKSA$ units. On the right side, the velocity dependence of $D_{\psi }^{\left( 0\right) }$ at $\xi_{0}=1$}{DradialJET}{7cm} %
%EndExpansion

The $3-D$ simulation is performed in the fully implicit mode with usual
parameters, i.e. $\Delta \overline{t}=10000,$ using the linearized
electron-electron Belaiev-Budker collision operator for momentum
conservation. While the time taken to calculate all matrix coefficients
associated to momentum dynamics is extremely fast, of the order of $120s%
\footnote{%
Absolute computer time values may depend significantly from the charge of
the computers. A reduction by $50\%$ is likely for a workstation with few
jobs.}$, since calculations may be performed in a vectorial form that is
very convenient for MatLab, coefficients for radial transport needs a much
longer computer time is this version of the code, typically an order of
magnitude more. This limitation results from the fact that loops have been
used in MatLab as a consequence of the complexity of the coefficients
arrangement in the matrix related to the change of the trapped/passing
boundary with radial position. In that case, a specific MEX-file written in
C or FORTRAN should considerably reduce time consumption, at a level
comparable to the one needed for calculating the coefficients which result
from momentum dynamics. 

In the case of full $3-D$ calculation, the Maxwellian solution is imposed at 
$i=1/2$. Though this condition is not necessary in principle, it turns out
to increase matrix conditionning, and reduce significantly the memory
required to store $LU$ matrices. However, no normalization of the density at
each time step is performed. Consequently, the drop tolerance for incomplete 
$LU$ matrix factorization may be increased from $10^{-4}$ to $10^{-3}$ as
compared to the case without radial transport, a important advantage since
it reduces considerably computer time consumption for performing the
incomplete factorization\footnote{%
It is interesting to note that the presence of a radial transport
contributes also to use larger drop tolerance parameter, by smoothing out
the momentum dynamics. It acts like a regularization process which increases
the matrix conditionning.}. On a Compaq workstation, the characteristic time
used for this calculation is around $1500s,$ and the memory required to
store $LU$ matrices is $214MBytes$ in the example here studied.

Once $LU$ matrices are calculated, the convergence is very fast. it is
achieved after $6$ iterations, corresponding to the minimum value required
for a correct convergence with the explicit momentum conservation term
(first Legendre truncated term) as discussed in Sec.\ref{fpinversionsection}%
. As shown in Fig.\ref{conservationprofJET}, the code is fully conservative,
within less than $1\%$.

%TCIMACRO{
%\TeXButton{Fig_conservationprofJET}{\fig{figures/conservationprofJET.eps}{Relative particle conservation of the drift kinetic code for the $3-D$ JET Lower Hybrid current drive simulation}
%{conservationprofJET}{7cm} }}
%BeginExpansion
\fig{figures/conservationprofJET.eps}{Relative particle conservation of the drift kinetic code for the $3-D$ JET Lower Hybrid current drive simulation}
{conservationprofJET}{7cm} %
%EndExpansion

While the absorbed RF power density profile given in Fig.\ref{powerprofJET}
shows a localized peak around $\rho =0.5,$ the current density profile is
much broader, as indicated in Fig.\ref{currentprofJET}. In the direction of
the plasma center, the broadening effect is weak, since collisional slowing
down prevails over radial transport when density is high. Towards plasma
edge, the effect of radial transport is much more visible, and for $\rho 
\geq 0.6,$ all the fast electrons are driven by this mechanism instead of
direct kinetic resonance. The effect of plasma magnetic equibrium on radial
transport is fully considered in the calculations. 

%TCIMACRO{
%\TeXButton{Fig_powerprofJET}{\fig{figures/powerprofJET.eps}{ Flux surface averaged power density profiles for collision, RF and Ohmic electric field absorption for the $3-D$
% JET Lower Hybrid current drive simulation. }
%{powerprofJET}{7cm} }}
%BeginExpansion
\fig{figures/powerprofJET.eps}{ Flux surface averaged power density profiles for collision, RF and Ohmic electric field absorption for the $3-D$
 JET Lower Hybrid current drive simulation. }
{powerprofJET}{7cm} %
%EndExpansion

%TCIMACRO{
%\TeXButton{Fig_currentprofJET}{\fig{figures/currentprofJET.eps}{ Flux surface averaged current density profiles for collision, RF and Ohmic electric field absorption for the $3-D$
% JET Lower Hybrid current drive simulation. }
%{currentprofJET}{7cm} }}
%BeginExpansion
\fig{figures/currentprofJET.eps}{ Flux surface averaged current density profiles for collision, RF and Ohmic electric field absorption for the $3-D$
 JET Lower Hybrid current drive simulation. }
{currentprofJET}{7cm} %
%EndExpansion

Comparison between distributions at radial positions where RF absorption
peaks (Figs.\ref{f2DpeakJET} and \ref{FparpeakJET}) and in the region where
radial transport predominates (Figs.\ref{f2DedgeJET} and \ref{FparedgeJET})
clearly shows the difference of dynamics in momentum space. While a plateau
is clearly formed at $\rho \simeq 0.36,$ the distribution exibits a bump at $%
\rho \simeq 0.78.$ There are nevertheless some reminiscence of the Lower
Hybrid quasilinear resonance boundaries. In all figures, the output of the
code is clean from numerical instabilities, even far from the resonance
domain, which confirms the robustness of the numerical scheme here employed
for $3-D$ calculations.

%TCIMACRO{
%\TeXButton{Fig_f2DpeakJET}{\fig{figures/f2DpeakJET.eps}{$2-D$ contour plot of the electron distribution function at $\rho\simeq 0.36$ for JET Lower Hybrid current drive}
%{f2DpeakJET}{7cm} }}
%BeginExpansion
\fig{figures/f2DpeakJET.eps}{$2-D$ contour plot of the electron distribution function at $\rho\simeq 0.36$ for JET Lower Hybrid current drive}
{f2DpeakJET}{7cm} %
%EndExpansion

%TCIMACRO{
%\TeXButton{Fig_FparpeakJET}{\fig{figures/FparpeakJET.eps}{Electron distribution function averaged over the perpendicular momentum direction at $\rho\simeq 0.36$
% for JET Lower Hybrid current drive.  The perpendicular and parallel temperatures are also shown}
%{FparpeakJET}{7cm} }}
%BeginExpansion
\fig{figures/FparpeakJET.eps}{Electron distribution function averaged over the perpendicular momentum direction at $\rho\simeq 0.36$
 for JET Lower Hybrid current drive.  The perpendicular and parallel temperatures are also shown}
{FparpeakJET}{7cm} %
%EndExpansion

%TCIMACRO{
%\TeXButton{Fig_f2DedgeJET}{\fig{figures/f2DedgeJET.eps}{$2-D$ contour plot of the electron distribution function at $\rho\simeq 0.78$ for JET Lower Hybrid current drive}
%{f2DedgeJET}{7cm} }}
%BeginExpansion
\fig{figures/f2DedgeJET.eps}{$2-D$ contour plot of the electron distribution function at $\rho\simeq 0.78$ for JET Lower Hybrid current drive}
{f2DedgeJET}{7cm} %
%EndExpansion

%TCIMACRO{
%\TeXButton{Fig_FparedgeJET}{\fig{figures/FparedgeJET.eps}{Electron distribution function averaged over the perpendicular momentum direction at $\rho\simeq 0.78$
% for JET Lower Hybrid current drive.  The perpendicular and parallel temperatures are also shown}
%{FparedgeJET}{7cm} }}
%BeginExpansion
\fig{figures/FparedgeJET.eps}{Electron distribution function averaged over the perpendicular momentum direction at $\rho\simeq 0.78$
 for JET Lower Hybrid current drive.  The perpendicular and parallel temperatures are also shown}
{FparedgeJET}{7cm} %
%EndExpansion

Finally, it is important to mention that the power density absorbed by
collisions is much broader than the RF contribution, as shown in Fig. \ref
{powerprofJET}. Since radial fast electron transport makes the current
response non-local, in that case the current drive efficiency may be only
defined as the ratio of the total current driven by the wave and the total
RF absorbed power in the plasma. 
